\section{An actual auxiliary cover compatible with the ordinary root construction}
\label{sec:cubic-actual-tame-cover}

The coefficient prime is \(p=269\).  We write \(\kappa\) for
the action on \(\mu_p\), so that
\(\kappa=\overline\epsilon^{-1}\) in the geometric Frobenius
convention of AC-G78.  Only the values \(1,-1\) of this
character enter the Kummer calculation.

\begin{lemma}[AC-G83-Kummer: adjusting a scalar without changing the ordinary branch]
\label{lem:G83-Kummer}
Let \(F/\mathbf Q\) be the degree-four G77 host, and let
\(\mathcal T\) be the continuous trace of the normalized
ordinary branch used in AC-G78.
Let \(E/\mathbf Q\) be a real quadratic field, split at \(p\).
Suppose there are \(h,t\in G_F\) with
\[
 \mathcal T(ha)=\sigma\mathcal T(a)\quad(a\in G_F),
 \quad \sigma\in\{1,-1\},\quad
 \epsilon(h)=\chi(h)=1,\quad \eta_E(h)=-1,
                                                        \tag{G83.1}
\]
and
\[
 \epsilon(t)=-(1+p),\quad \chi(t)=1,\quad\eta_E(t)=1.
                                                        \tag{G83.2}
\]
Let \(u\) be the square of a fundamental unit of \(E\), and set
\[
 M=E(\mu_p,u^{1/p}),\qquad
 A=\operatorname{Gal}(M/E(\mu_p)).
\]
Then one can replace \(h\) by a conjugate \(h'=nhn^{-1}\),
with \(n\in G_F\), such that
\[
                    (h't)^2|_M=1.                \tag{G83.3}
\]
All scalar trace identities (G83.1) and all abelian
character values of \(h\) are unchanged.
\end{lemma}

\begin{proof}
The unit \(u\) has norm one, so its nontrivial quadratic
conjugate is \(u^{-1}\).  Therefore \(M/\mathbf Q\) is Galois.
The Kummer group \(A\) is trivial or cyclic of order \(p\).
On it an element \(b\in G_{\mathbf Q}\) acts by
\[
                       \eta_E(b)\kappa(b).         \tag{G83.4}
\]
Indeed, \(b\) acts on the class of \(u\) by the exponent
\(\eta_E(b)\), and on the root-of-unity multiplier of
\(u^{1/p}\) by \(\kappa(b)\).
In particular both \(h\) and \(t\) act on \(A\) by inversion.
Their product \(g=ht\) acts trivially on \(A\).
Its image on \(E(\mu_p)\) has order two, so \(g^2\in A\).

If \(A=1\), nothing is required.  Otherwise the image of
\(G_F\) in \(\operatorname{Gal}(M/\mathbf Q)\) contains \(A\).
To see this without assuming disjointness, its index is
\([F\cap M:\mathbf Q]\), a divisor of \(4\).
The unique normal subgroup \(A\) of order \(p\) must therefore
lie in that image: its intersection has index dividing both
\(p\) and this prime-to-\(p\) index.
Choose any \(a_0\in A\), and lift it to \(n\in G_F\).
In \(\operatorname{Gal}(M/\mathbf Q)\) we then have
\[
 nhn^{-1}=a_0^2h,\qquad
 (nhn^{-1}t)^2=a_0^4g^2.                           \tag{G83.5}
\]
The first equality uses inversion by \(h\), and the second
uses the trivial action of \(g\).
Fourth powers are an automorphism of \(A\).
Choosing \(a_0^4=(g^2)^{-1}\) proves (G83.3).

Every abelian character is constant under conjugation.
The trace calculation uses its cyclicity:
\[
 \mathcal T(nhn^{-1}a)
    =\mathcal T(hn^{-1}an)
    =\sigma\mathcal T(n^{-1}an)
    =\sigma\mathcal T(a).
\]
This is valid for the actual continuous branch trace, not
only for its specialization or residual trace.
\end{proof}

\begin{proposition}[AC-G83-ray: a totally real dihedral ray cover with controlled degree]
\label{prop:G83-ray}
Let \(E\) be real quadratic and let \(q\ne p\) be a rational
prime inert in \(E\), with
\[
 v_p(q+1)=1,\qquad
 u\bmod q\in(\mathbf F_{q^2}^{\times})^p,            \tag{G83.6}
\]
where \(u\) is the square of a fundamental unit of \(E\).
Then there is a totally real extension \(T/E\), cyclic of
degree \(p^s\), such that \(T/\mathbf Q\) is dihedral, \(T/E\)
is unramified away from the unique prime above \(q\), and its
inertia group there has order exactly \(p\).  Moreover
\[
 1\le s\le1+v_p(h_E),\qquad
 p^s\le p h_E\le p B_E(1+\log B_E),\qquad
 B_E=\frac{\sqrt{\operatorname{Disc}E}}2.            \tag{G83.7}
\]
\end{proposition}

\begin{proof}
Since \(q\equiv-1\pmod p\), the \(p\)-primary part of
\(\mathbf F_{q^2}^{\times}\) is cyclic of order \(p\).
Condition (G83.6) says that \(u\) has trivial projection to
this part.  The same holds for a fundamental unit, because
squaring is invertible on a \(p\)-group, and for \(-1\).
Thus every global unit has trivial \(p\)-primary residue.

For the modulus consisting of the finite prime
\(\mathfrak q=q\mathcal O_E\), with no infinite primes, the
ray-class exact sequence gives
\[
 0\longrightarrow C_p\longrightarrow
       \operatorname{Cl}_{\mathfrak q}(E)(p)
       \longrightarrow\operatorname{Cl}(E)(p)
       \longrightarrow0.                         \tag{G83.8}
\]
This is the ordinary ray group, not the narrow group.
The nontrivial automorphism \(c\) of \(E/\mathbf Q\) acts on
the left term by \(x\mapsto x^q=x^{-1}\).
It acts on the ordinary ideal class group by inversion,
since an ideal times its conjugate is a rational principal
ideal.  The idempotent \((1+c)/2\) is exact on finite abelian
\(p\)-groups.  Applying it to (G83.8) shows that \(c\) acts
by inversion on the entire middle group.

Choose a faithful character of its subgroup \(C_p\) into
\(\mathbf Q_p/\mathbf Z_p\).
Divisibility of that target extends the character to the
whole finite abelian \(p\)-group.  Its image is cyclic of
order \(p^s\), and its kernel is \(c\)-stable because \(c\)
acts by inversion.
Global class field theory now gives the corresponding
cyclic extension \(T/E\).  The modulus has no real places,
so \(T\) is totally real.  It has no finite ramification
away from \(\mathfrak q\), and the local inertia image is
exactly \(C_p\), because the chosen character is faithful
there.  Its kernel is \(c\)-stable, so \(T/\mathbf Q\) is
Galois and \(c\) acts on its cyclic subgroup by inversion.
A lift of \(c\) has square fixed by inversion, hence square
one in the odd-order cyclic group.  The full Galois group
is therefore dihedral.

The exact sequence has order \(p\,|\operatorname{Cl}(E)(p)|\),
which proves the first two bounds in (G83.7).
For the last bound, the real quadratic Minkowski bound
represents every ideal class by an integral ideal of norm
at most \(B_E\).  For a quadratic field the number \(a_E(n)\)
of integral ideals of norm \(n\) satisfies \(a_E(n)\le d(n)\):
at a split prime the local count is \(k+1\), at an inert
prime it is at most one, and at a ramified prime it is one.
Consequently
\[
 h_E\le\sum_{n\le B_E}d(n)
      =\sum_{ab\le B_E}1
      \le B_E\sum_{a\le B_E}\frac1a
      \le B_E(1+\log B_E).
\]
Here \(B_E>1\) for a real quadratic field.
The exact ray sequence is Milne, \emph{Class Field Theory}
v4.03, Chapter~V, Theorem~1.7, printed pp.~150--151;
the existence and conductor assertions are V, Theorems~3.5--3.6
and Remark~3.8, printed p.~158.
Minkowski's input is Milne, \emph{Algebraic Number Theory}
v3.08, Theorem~4.3, printed p.~70.
Global class field theory and Minkowski's theorem are
retained standard inputs; the character and degree
arguments above are proved here.
\end{proof}

\begin{theorem}[AC-G83: the auxiliary type can be killed on an actual compatible CM cover]
\label{thm:G83}
Consider either of the following actual ordinary constructions.

In the odd branch, take the inherited odd-\(e\) input of
AC-G78--AC-G80 on its G77 host \(F\).
Choose, before \(q\), a fresh real quadratic field
\(E=\mathbf Q(\sqrt{r_2})\), where \(r_2\) is a prime
congruent to \(1\) modulo \(8p\) and every fixed old prime.

In the positive-sign branch of AC-G81, keep its fields
\(D=\mathbf Q(\sqrt{-r_0})\), \(K'=\mathbf Q(\sqrt{-r_1})\),
\(F=KK'\), \(L=FD\), and put
\[
                         E=\mathbf Q(\sqrt{r_0r_1})
                         \subset L^+.
\]
Thus this branch retains the separated residual condition
\(\overline\gamma^2|_{G_K}\ne1\).
In both branches the actual construction uses the simultaneous
projector \(e_{\rm aux}e_{\rm old}e_{\delta_e}\) from the start.
For AC-G81 its old-type list includes the exact scalar
\(\delta_0\) on both endpoint lines at \(r_0\), with the
middle line trivial.  The proof below verifies availability
of these projectors for the twisted input.

The prime \(q\) in the actual ordinary root construction
can be chosen so that a cover \(T/E\) of
Proposition~\ref{prop:G83-ray} exists.
Put \(H=FT\) in the odd branch and \(H=LT\) in the
positive-sign branch.
Then
\[
 [H:\mathbf Q]=8p^s\le8p h_E,\qquad
 H/H^+\ \hbox{is unramified at all finite places}.   \tag{G83.9}
\]
After soluble automorphic base change and unitary descent
there is an actual ordinary rank-three automorphic
representation on the definite unitary group for \(H/H^+\)
with the original symmetric weight and original cubic
residual parameter restricted to \(G_H\).
Its normalized Galois representation is absolutely
irreducible.  It is unramified at every place above the
new auxiliary \(q\); it retains the exact old away inertia
types when the simultaneous AC-G80 projectors are imposed,
and retains the exact ordered coefficient-prime flag.
In the AC-G81 branch the added scalar twist is also killed.

Writing \(\operatorname{rd}\) for root discriminant, one has
\[
 \operatorname{rd}(H)=
 \begin{cases}
   \operatorname{rd}(F)\,r_2^{1/2}\,
           q^{\,1-1/p},&\text{odd branch},\\
   \operatorname{rd}(L)\,q^{\,1-1/p},
           &\text{positive-sign branch}.
 \end{cases}                                      \tag{G83.10}
\]
This verifies existence of an actual cover for the selected
\(q\), rather than assuming the earlier AC-G73 cover
hypotheses.  It gives no bound on \(q\), no numerical frame
or counting-budget conclusion, and no per-curve
\(J\)-containment.
\end{theorem}

\begin{proof}
\emph{The two scalar choices.}
For the odd branch, the ordinary Hida trace is unramified
at the fresh \(r_2\).  Choose \(h\in I_{r_2}\) with
\(\eta_E(h)=-1\).
Because all fixed data and the whole branch are unramified
there, \(h\in G_F\),
\(\epsilon(h)=\chi(h)=1\), and
\(\mathcal T(ha)=\mathcal T(a)\) for every \(a\).
The choice of congruence for \(r_2\) makes \(E\) split at \(p\).
The ordinary inertia element \(t\) in AC-G78 has
\(\epsilon(t)=-(1+p)\), \(\chi(t)=1\), and its trace on
the branch is the required negative function of half-weight.
Since \(E\) is unramified at \(p\), \(\eta_E(t)=1\).
Thus the product \(ht\) gives exactly the old analytic
target.

For the positive-sign branch, take the scalar \(h\) of
AC-G78 applied to the fresh twist of AC-G81.
It has \(\mathcal T(ha)=-\mathcal T(a)\),
\(\epsilon(h)=\chi(h)=1\),
\(\delta_0(h)=-1\), and \(\delta_1(h)=1\), where
\(\delta_i=\eta_{-r_i}\).
The last two values hold for the original fresh inertia
witness and remain unchanged under its odd \(p\)-power
limit used to construct the scalar on the complete branch.
Since \(\eta_E=\delta_0\delta_1\), this gives
\(\eta_E(h)=-1\).
Both quadratic characters are trivial on inertia at \(p\),
and both split at \(p\); the same \(t\) has
\(\eta_E(t)=1\).
Its ordinary inertia sign is positive, so \(ht\) again
gives precisely the required negative analytic trace.

Apply Lemma~\ref{lem:G83-Kummer} in each case.
Conjugating \(h\) changes none of these trace or character
identities.  Therefore the entire Weierstrass argument in
AC-G78, including its treatment of a ramified branch,
remains unchanged.  In the finite quotient used for
Chebotarev add the action on \(M\), and also prescribe
\(\epsilon\bmod p^2\).
Chebotarev still approximates the actual element \(h't\);
there is no independently imposed incompatible Frobenius
condition.
Its conjugacy class is enough: the trace, abelian values
and square-one condition are invariant under conjugation.

The resulting rational prime \(q\) is split completely
in \(F\), has \(\chi(\operatorname{Frob}_q)=1\),
is inert in \(E\), and satisfies
\[
 q\equiv-1\pmod p,\qquad v_p(q+1)=1.
\]
The last assertion follows from
\(\epsilon(\operatorname{Frob}_q)=q^{-1}
 \equiv-(1+p)\pmod{p^2}\).
Its Frobenius on \(M\) has square one.
At the inert prime of \(E\), the residue field is
\(\mathbf F_{q^2}\), which contains \(\mu_p\).
The square-one condition says that this prime splits
in the Kummer extension over \(E(\mu_p)\).
Reducing \(u^{1/p}\) therefore shows
\(u\bmod q\in(\mathbf F_{q^2}^{\times})^p\).
All fields are fixed before this Chebotarev choice, and
\(q\) is chosen away from their ramified primes.
Proposition~\ref{prop:G83-ray} now constructs \(T\).

\emph{The ordinary seed is produced before the cover.}
The modified finite quotient changes neither the exact
analytic trace nor the actual Hecke-support argument.
Thus the four-test Steinberg conclusion, projective
finite-type transfer, global irreducibility and
coefficient-prime flag of AC-G78--AC-G81 all apply to
this chosen \(q\).

The exact old away types can be selected simultaneously.
For the odd branch this is AC-G80 as stated.
For the positive-sign branch its proof is unchanged:
the fresh twist has at \(r_0\) the scalar finite inertia
\(\delta_0\) on its two endpoint lines, of order two,
and at every other old away prime it has exactly the
previous finite inertia, up to the unramified twist.
The same finite-character idempotents and Saha
specialization argument apply.
The compact characters still have order prime to \(269\).
The projectors \(e_{\rm aux}e_{\rm old}e_{\delta_e}\)
are imposed from the start, and commute with every
trace, four-test map and \(q\)-level finite-group action.
The additional Kummer quotient in Chebotarev does not
alter these modules or maps.
This supplies the simultaneous old types without
assuming an old \(\mathrm{Dat}\) for the twisted character.

\emph{Fields and ramification.}
The odd-order dihedral extension \(T/\mathbf Q\) has
exactly one quadratic subfield, namely \(E\).
In the odd branch the fresh \(r_2\) gives \(E\not\subset F\).
Since \([F:\mathbf Q]=4\), it follows that \(F\cap T=\mathbf Q\).
Hence \([FT:\mathbf Q]=8p^s\), with the tower
\(F\subset FE\subset FT\) having a quadratic first step
and a cyclic \(p^s\) second step.
In the positive-sign branch, \(E\subset L\), and degree
considerations give \(L\cap T=E\).
Thus \([LT:\mathbf Q]=8p^s\), cyclic of degree \(p^s\)
over \(L\).
Since \(T\) is totally real, the composita are CM, with
maximal real fields \(F^+T\) and \(L^+T\), respectively.
Their quadratic CM extensions are base changes of
\(F/F^+\) and \(L/L^+\), so are unramified at all finite
places and split above \(p\).

At \(q\) the relative ramification index from the ray
cover is exactly \(p\).
In the positive-sign case \(q\) is unramified quadratic
in \(L/F\); in the odd case \(E\) supplies that unramified
quadratic extension at \(q\).
Consequently the old tame two-dimensional local
parameter first splits into its order-\(p\) inertia
characters, and these characters become trivial after
the ramification-index-\(p\) extension.
The unramified line remains unramified, and \(N=0\)
remains zero.  Thus the whole parameter is unramified
above \(q\).
There is no new finite ramification of \(T/E\) elsewhere.
In the odd branch \(E/\mathbf Q\) ramifies at \(r_2\),
but the automorphic parameter was unramified there.
In the positive-sign branch the character \(\delta_0\)
is killed on \(L\); the selected scalar old type at
\(r_0\) is therefore unramified.
The remaining prescribed old inertia is the restriction
of the original one.

The cyclic degree-\(p\) and quadratic steps preserve
absolute irreducibility of the three-dimensional
representation.  The proof of (G81.8) works for any
cyclic prime index \(\ell\): restriction is semisimple
and its endomorphism algebra is the sum of
\(\operatorname{Hom}_G(W,W\otimes\nu)\) over characters
of that cyclic quotient.  A nontrivial self-twist
would imply \(\nu^3=1\) by determinants.
For \(\ell=2\) or \(269\), this is impossible.
Iterating proves irreducibility over \(H\).
Residual separation survives as well.  In the odd branch,
if \(\overline\gamma^2|_{G_F}\) became trivial on \(G_{FE}\),
it would equal the nontrivial quadratic character of
\(FE/F\); the former is unramified at the fresh \(r_2\),
whereas the latter is ramified there.
In the positive-sign branch separation on \(G_L\) is the
fresh-\(r_0,r_1\) argument of AC-G81.
Subsequent extensions of \(p\)-power degree cannot kill a
nontrivial residual character of order prime to \(p\).
The cubic character also retains order three throughout
these degree-\(2\) and degree-\(269\) steps.

Clozel--Thorne I, Theorem~2.7(1), gives cuspidal soluble
automorphic base change along this tower.
The source ramification remains at relative split
places; new field ramification cannot ramify an
unramified representation.
Thorne's Theorem~4.4 therefore gives the actual definite
unitary representation on \(H/H^+\).
The coefficient-prime extensions are unramified:
the fresh real quadratic is split there and \(T/E\)
has conductor only at \(q\).
The actual G79 flag restricts to these extensions,
retaining its ordered characters and parallel Hodge
weights.  Ordinary slopes and the sums of parallel
weights are multiplied by the same residue degree,
so ordinarity is retained; equivalently use
Clozel--Thorne I, Theorem~2.7(3).
No residual-ordinary-to-characteristic-zero inference
is used.

Finally, a tame Galois extension with ramification
index \(p\) has different exponent \(p-1\) locally.
All \(q\)-places have that index, so the relative
contribution to root discriminant is
\(q^{1-1/p}\).
In the odd branch \(r_2\) is fresh for \(F\);
the quadratic compositum \(FE/F\) contributes
\(r_2^{1/2}\), while in the positive-sign branch
the starting field is already \(L\).
This proves (G83.10).
The degree estimate is (G83.7).
The construction supplies the cover, but does not
bound the nonabelian Chebotarev approximation prime
or verify any numerical frame or actual \(J\)-support
condition.
\end{proof}
